\section{Conclusion}\label{sec:conclusion}

This paper asks when heterogeneous AV companies should coordinate their
routing and when partial independence should be retained. Companies experience
the same physical road costs but can manage their fleets under different
objectives. A coalition aggregates those objectives and internalizes selected
cross-member effects while preserving member OD obligations; HDVs continue to
choose routes through a nonatomic Wardrop equilibrium. The resulting
organizational problem concerns the continuation congestion induced by a
specified coalition, not the ownership shares alone and not the private
stability of coalition formation.

The theory connects objective heterogeneity to a usable coalition-design
criterion. On a fixed affine support, the activated signature retains the
OD-feasible coalition response that survives HDV reallocation. Its projection
onto marginal system cost yields a general-network merge-direction score. The
canonical bottleneck exposes three regimes: HDV screening can make organization
neutral; a boundary regime favors changes that lower effective response; and
an interior-response regime can make partial coordination better than both
grand and singleton organization. Adding governance burden gives closed-form
pairwise thresholds and a lower envelope of grand, partial, and fragmented
policy intervals.

The certified experiments test both the regime and the selection rule. Sioux
Falls moves from a grand-coalition optimum at moderate penetration to partial
coordination at high penetration. Independent Braess and directed-grid scans
show the same qualitative transition, with topology-specific boundaries. A
one-step merge policy reaches an audit optimum in $33$ of $36$ total states and
has mean regret $0.332$ points on the $32$ cross-network states; two-step
lookahead reaches an audit optimum in all $36$. The latter result is evidence
within the declared four-company designs rather than a universal guarantee.
Objective-composition, company-count, spatial-exposure, and sensitivity tests
explain why the policy cannot be replaced by HHI or coalition count.

The practical conclusion is conditional but specific. A decision maker should
estimate or simulate the activated response, use its system-cost projection to
screen candidate merges, search a local merge neighborhood, and account for
the cost of cross-company governance. Centralization is justified when that
process moves the network toward its response target and the congestion gain
covers the coordination burden. Partial independence is justified when further
coupling overshoots the target or costs more than it saves. Extending this
continuation design to empirically calibrated objectives and endogenous
coalition formation is a separate research step.
